\documentclass[12pt,a4paper,oneside]{book}

\usepackage{fontspec}
\usepackage{polyglossia}
\setmainlanguage{russian}
\setotherlanguage{english}
\setmainfont{Times New Roman}

\usepackage{amsmath,amssymb,amsthm,mathtools}
\usepackage{geometry}
\usepackage{enumitem}
\usepackage{hyperref}
\usepackage{microtype}

\geometry{margin=2.5cm}
\setlength{\parindent}{1.25cm}
\setlength{\parskip}{0.35em}
\sloppy

\hypersetup{
    colorlinks=true,
    linkcolor=black,
    citecolor=black,
    urlcolor=blue,
    pdftitle={Теория Единого Целого},
    pdfauthor={}
}

\theoremstyle{definition}
\newtheorem{axiom}{Аксиома}[chapter]
\newtheorem{definition}{Определение}[chapter]
\newtheorem{postulate}{Постулат}[chapter]

\theoremstyle{plain}
\newtheorem{theorem}{Теорема}[chapter]
\newtheorem{lemma}{Лемма}[chapter]
\newtheorem{corollary}{Следствие}[chapter]
\newtheorem{proposition}{Предложение}[chapter]

\theoremstyle{remark}
\newtheorem{remark}{Замечание}[chapter]
\newtheorem{example}{Пример}[chapter]

\newcommand{\Whole}{U}
\newcommand{\Parts}{\mathcal A}
\newcommand{\SubParts}{\mathcal B}
\newcommand{\Rel}{R}
\newcommand{\State}{S}
\newcommand{\States}{\mathcal S}
\newcommand{\Values}{\mathcal V}
\newcommand{\DistVals}{\mathcal D}
\newcommand{\Geom}{\mathcal G}
\newcommand{\Space}{\mathsf{Space}}
\newcommand{\Time}{\mathsf{Time}}
\newcommand{\Physics}{\mathsf{Physics}}
\newcommand{\Motion}{\mathsf{Motion}}
\newcommand{\Observer}{\mathsf{Obs}}
\newcommand{\Struct}{\operatorname{Struct}}
\newcommand{\Aut}{\operatorname{Aut}}
\newcommand{\Inv}{\operatorname{Inv}}
\newcommand{\Diff}{\operatorname{Diff}}
\newcommand{\Card}{\operatorname{card}}
\newcommand{\Id}{\operatorname{id}}

\title{\Huge\textbf{Теория Единого Целого}\\[0.8em]
\Large Аксиоматика, теория отношений и выведение пространства-времени}
\author{}
\date{}

\begin{document}

\frontmatter
\maketitle

\chapter*{Аннотация}
\addcontentsline{toc}{chapter}{Аннотация}

Настоящий текст представляет собой формальную монографическую версию Теории Единого Целого. Основная цель теории состоит в построении языка, в котором пространство, время, движение, расстояние, масса и энергия не вводятся как первичные понятия, а возникают как производные структуры из более фундаментального уровня: Единого Целого, его частей и отношений между ними.

Исходная идея может быть выражена кратко:
\[
\boxed{\text{не отношения существуют в пространстве и времени,}}
\]
\[
\boxed{\text{а пространство и время возникают из отношений.}}
\]

В монографии последовательно вводятся первичные символы, аксиомы, определения, первые теоремы, теория состояний, реляционная геометрия, реляционное время, наблюдатели, движение и программа физической интерпретации.

\tableofcontents

\mainmatter

\chapter{Введение}

\section{Мотивировка}

Большинство физических теорий начинают построение с уже заданных понятий пространства и времени. В них располагаются объекты, происходят события и описывается движение. Такая схема имеет вид
\[
\text{пространство} + \text{время} + \text{объекты} + \text{законы}.
\]

Теория Единого Целого предлагает противоположный порядок. Она начинает не с пространства и времени, а с более абстрактных понятий:
\[
\Whole,\quad \Parts,\quad \Rel,\quad \State.
\]

Здесь \(\Whole\) обозначает Единое Целое, \(\Parts\) --- класс его частей, \(\Rel\) --- отношения между частями, а \(\State\) --- состояние системы отношений. Пространство и время должны быть выведены позднее.

\section{Основная идея}

Фундаментальная последовательность теории имеет вид
\[
\Whole
\longrightarrow
\Parts
\longrightarrow
\Rel
\longrightarrow
\State
\longrightarrow
\Space
\longrightarrow
\Time
\longrightarrow
\Physics.
\]

Расшифровка этой схемы такова:
\[
\begin{array}{rcl}
\Whole &\Rightarrow& \text{Единое Целое},\\
\Parts &\Rightarrow& \text{части},\\
\Rel &\Rightarrow& \text{отношения между частями},\\
\State &\Rightarrow& \text{состояния отношений},\\
\Space &\Rightarrow& \text{пространство как структура отношений},\\
\Time &\Rightarrow& \text{время как изменение отношений},\\
\Physics &\Rightarrow& \text{физика как теория устойчивых реляционных структур}.
\end{array}
\]

\section{Методологическое требование}

В данной теории запрещается использовать пространство, время, координаты, расстояние, скорость, массу и энергию в качестве исходных понятий. Эти понятия могут появиться только после того, как будут определены отношения и состояния.

Это требование можно записать как принцип вывода:
\[
\Rel \prec \Space,\quad \Rel \prec \Time,
\]
где \(\prec\) означает отношение концептуальной первичности.

\chapter{Формальный язык}

\section{Первичные символы}

Вводятся следующие первичные неопределяемые символы:
\[
\Whole,\quad \Parts,\quad \Rel,\quad \State.
\]

Их интуитивный смысл:
\[
\begin{array}{rcl}
\Whole &:& \text{Единое Целое},\\
\Parts &:& \text{класс частей Единого Целого},\\
\Rel &:& \text{отношение между частями},\\
\State &:& \text{состояние системы отношений}.
\end{array}
\]

\section{Запрещённые первичные понятия}

На исходном уровне языка не допускаются как первичные:
\[
 x,\quad t,\quad d,\quad v,\quad m,\quad E.
\]

Иначе говоря, координата, время, расстояние, скорость, масса и энергия не принадлежат начальному словарю теории.

\begin{remark}
Это не означает, что данные величины не существуют. Это означает только, что они не являются фундаментальными. Они должны быть получены как производные от \(\Rel\) и \(\State\).
\end{remark}

\section{Класс частей}

Пусть \(\Parts\) обозначает класс всех частей Единого Целого:
\[
\Parts = \{A_i \mid A_i \subset \Whole,\ A_i \neq \Whole\}.
\]

Элементы \(A_i \in \Parts\) называются частями.

\chapter{Аксиоматика Единого Целого}

\begin{axiom}[Единственность Целого]
Существует единственный объект \(\Whole\), называемый Единым Целым:
\[
\exists!\,\Whole.
\]
\end{axiom}

\begin{axiom}[Всеобщая принадлежность]
Для всякого существующего объекта \(X\) выполнено
\[
X \subseteq \Whole.
\]
\end{axiom}

\begin{axiom}[Собственность частей]
Всякая часть является собственной частью Единого Целого:
\[
\forall A \in \Parts:\quad A \subset \Whole,
\]
и
\[
\forall A \in \Parts:\quad A \neq \Whole.
\]
\end{axiom}

\begin{axiom}[Отсутствие внешнего]
Не существует объекта, внешнего по отношению к \(\Whole\):
\[
\nexists X:\quad X \not\subseteq \Whole.
\]
\end{axiom}

\begin{axiom}[Реляционность свойств]
Никакое свойство не принадлежит части само по себе. Для всякой части \(A_i \in \Parts\) её свойства определяются только через отношения с другими частями:
\[
P(A_i) := \{\Rel(A_i,A_j) \mid A_j \in \Parts,\ A_j \neq A_i\}.
\]
\end{axiom}

\begin{axiom}[Первичность отношений]
Отношение между частями является более первичным, чем метрические и физические величины. В исходном языке не существуют самостоятельные функции
\[
d(A_i,A_j),\quad t,\quad v(A_i,A_j),\quad m(A_i),\quad E(A_i).
\]
Они могут быть введены только как производные от \(\Rel\) и \(\State\).
\end{axiom}

\section{Непосредственные следствия аксиом}

\begin{theorem}[Невозможность внешней системы координат]
В Теории Единого Целого не существует внешней абсолютной системы координат.
\end{theorem}

\begin{proof}
Внешняя система координат требует внешней области отсчёта. По Аксиоме отсутствия внешнего не существует объекта или области вне \(\Whole\). Следовательно, внешняя система координат не определена.
\end{proof}

\begin{theorem}[Невозможность абсолютного внешнего времени]
В Теории Единого Целого не существует внешнего абсолютного времени.
\end{theorem}

\begin{proof}
Абсолютное внешнее время предполагало бы параметр, заданный независимо от всех отношений внутри \(\Whole\). Но всё существующее принадлежит \(\Whole\), а свойства частей определяются только через отношения. Следовательно, внешний временной параметр не принадлежит исходному языку теории.
\end{proof}

\begin{theorem}[Отсутствие абсолютного центра]
В Едином Целом не существует абсолютного центра.
\end{theorem}

\begin{proof}
Абсолютный центр может быть определён только относительно заранее заданной внешней координатной структуры. Но такая структура невозможна по Аксиоме отсутствия внешнего. Следовательно, абсолютный центр не определён.
\end{proof}

\chapter{Теория отношений}

\section{Отношение как отображение}

\begin{definition}[Отношение]
Отношением называется отображение
\[
\Rel:\Parts \times \Parts \to \Values,
\]
где \(\Values\) --- множество возможных значений отношений.
\end{definition}

Для двух частей \(A_i,A_j \in \Parts\) значение
\[
\Rel(A_i,A_j) \in \Values
\]
называется отношением между \(A_i\) и \(A_j\).

\section{Виды отношений}

Отношение \(\Rel\) может обладать различными свойствами.

\begin{definition}[Симметричное отношение]
Отношение называется симметричным, если
\[
\Rel(A_i,A_j)=\Rel(A_j,A_i)
\]
для всех \(A_i,A_j \in \Parts\).
\end{definition}

\begin{definition}[Направленное отношение]
Отношение называется направленным, если вообще говоря
\[
\Rel(A_i,A_j)\neq \Rel(A_j,A_i).
\]
\end{definition}

\begin{definition}[Рефлексивное отношение]
Отношение называется рефлексивным, если для всякой части \(A_i\) определено
\[
\Rel(A_i,A_i).
\]
\end{definition}

\begin{remark}
В базовой версии ТЕЦ отношения между различными частями имеют приоритет над самотношениями. Поэтому часто рассматривается множество пар
\[
\Parts^{(2)}=\{(A_i,A_j)\in \Parts\times\Parts\mid A_i\neq A_j\}.
\]
\end{remark}

\section{Система отношений}

\begin{definition}[Система частей]
Системой частей называется непустое подмножество
\[
\SubParts \subseteq \Parts.
\]
\end{definition}

\begin{definition}[Реляционная структура]
Реляционной структурой системы \(\SubParts\) называется пара
\[
\mathfrak R(\SubParts) := (\SubParts,\Rel|_{\SubParts\times\SubParts}).
\]
\end{definition}

\begin{definition}[Состояние]
Состоянием системы \(\SubParts\) называется совокупность всех отношений между её различными частями:
\[
\State(\SubParts):=\{\Rel(A_i,A_j)\mid A_i,A_j\in\SubParts,\ i\neq j\}.
\]
\end{definition}

\section{Минимальные системы}

\begin{theorem}[Одна часть не образует пространства]
Если
\[
\SubParts=\{A_1\},
\]
то в \(\SubParts\) отсутствуют отношения между различными частями, а значит не возникает пространство.
\end{theorem}

\begin{proof}
Для существования отношения между различными частями нужны \(A_i,A_j\in\SubParts\), где \(i\neq j\). Если \(\SubParts\) содержит только \(A_1\), такой пары нет. Следовательно,
\[
\State(\SubParts)=\varnothing.
\]
По Аксиоме реляционности свойств без отношений не возникает структура различимости. Пространство как структура различимости не возникает.
\end{proof}

\begin{theorem}[Две части задают минимальное различие]
Если
\[
\SubParts=\{A_1,A_2\},\quad A_1\neq A_2,
\]
то возможно минимальное отношение
\[
\Rel(A_1,A_2).
\]
\end{theorem}

\begin{proof}
Пара \((A_1,A_2)\) принадлежит \(\SubParts\times\SubParts\), и её элементы различны. По определению отношения существует значение \(\Rel(A_1,A_2)\in\Values\), если отношение задано на этой паре. Это и есть минимальная структура различия.
\end{proof}

\chapter{Состояния и изменение}

\section{Множество состояний}

Пусть \(\States(\SubParts)\) обозначает множество возможных состояний системы \(\SubParts\):
\[
\States(\SubParts)=\{\State_k(\SubParts)\mid k\in K\}.
\]

Индекс \(k\) здесь не является временем. Он лишь различает возможные состояния.

\begin{definition}[Различие состояний]
Два состояния различны, если
\[
\State_1(\SubParts)\neq \State_2(\SubParts).
\]
\end{definition}

\begin{definition}[Изменение]
Изменением называется упорядоченная пара различных состояний:
\[
\Delta\State := (\State_1,\State_2),\quad \State_1\neq\State_2.
\]
\end{definition}

\section{Функция различия}

\begin{definition}[Мера различия состояний]
Мерой различия состояний называется отображение
\[
T:\States(\SubParts)\times\States(\SubParts)\to \mathbb R_{\ge 0},
\]
такое что
\[
T(\State_1,\State_2)=0 \iff \State_1=\State_2.
\]
\end{definition}

\begin{definition}[Реляционное время]
Реляционным временем между двумя состояниями называется величина
\[
t(\State_1,\State_2):=T(\State_1,\State_2).
\]
\end{definition}

\begin{theorem}[Время возникает из изменения]
Если \(\State_1\neq\State_2\), то возникает ненулевая реляционная временная величина.
\end{theorem}

\begin{proof}
По определению меры различия состояний
\[
T(\State_1,\State_2)=0 \iff \State_1=\State_2.
\]
Если \(\State_1\neq\State_2\), то
\[
T(\State_1,\State_2)>0.
\]
Следовательно, реляционное время возникает как мера различия состояний.
\end{proof}

\begin{theorem}[При неизменности отношений время не возникает]
Если
\[
\Rel_1(A_i,A_j)=\Rel_2(A_i,A_j)
\]
для всех различных \(A_i,A_j\in\SubParts\), то
\[
t(\State_1,\State_2)=0.
\]
\end{theorem}

\begin{proof}
Если все отношения совпадают, то совпадают и состояния:
\[
\State_1(\SubParts)=\State_2(\SubParts).
\]
Следовательно, по определению \(T\),
\[
t(\State_1,\State_2)=T(\State_1,\State_2)=0.
\]
\end{proof}

\chapter{Выведение пространства}

\section{Пространство как структура отношений}

\begin{definition}[Реляционное пространство]
Реляционным пространством системы \(\SubParts\) называется структура
\[
\Space(\SubParts):=\Struct(\State(\SubParts)).
\]
\end{definition}

То есть пространство определяется не как вместилище частей, а как структурная форма их отношений.

\section{Реляционное расстояние}

\begin{definition}[Реляционное расстояние]
Пусть задано отображение
\[
F:\Values\to\DistVals.
\]
Тогда реляционным расстоянием между частями \(A_i\) и \(A_j\) называется
\[
d(A_i,A_j):=F(\Rel(A_i,A_j)).
\]
\end{definition}

\begin{theorem}[Расстояние не первично]
Расстояние между частями не является первичным понятием ТЕЦ.
\end{theorem}

\begin{proof}
По определению расстояние выражается через отношение:
\[
d(A_i,A_j)=F(\Rel(A_i,A_j)).
\]
Следовательно, для определения \(d\) необходимо предварительно иметь \(\Rel\). Значит, \(d\) производно от \(\Rel\).
\end{proof}

\section{Метрические условия}

Если множество \(\DistVals\subseteq\mathbb R_{\ge0}\), то можно потребовать метрические условия:
\[
d(A_i,A_j)\ge0,
\]
\[
d(A_i,A_j)=0\iff A_i=A_j,
\]
\[
d(A_i,A_j)=d(A_j,A_i),
\]
\[
d(A_i,A_k)\le d(A_i,A_j)+d(A_j,A_k).
\]

\begin{definition}[Метризуемая реляционная система]
Система \(\SubParts\) называется метризуемой, если существует отображение \(F\), такое что функция
\[
d(A_i,A_j)=F(\Rel(A_i,A_j))
\]
удовлетворяет аксиомам метрики.
\end{definition}

\begin{remark}
Не всякая система отношений обязана быть метризуемой. Это важно: пространство в привычном метрическом смысле является частным случаем более общей реляционной структуры.
\end{remark}

\section{Геометрия как инвариант отношений}

\begin{definition}[Автоморфизм реляционной структуры]
Автоморфизмом реляционной структуры \(\mathfrak R(\SubParts)\) называется биекция
\[
\phi:\SubParts\to\SubParts,
\]
сохраняющая отношения:
\[
\Rel(A_i,A_j)=\Rel(\phi(A_i),\phi(A_j)).
\]
\end{definition}

\begin{definition}[Геометрический инвариант]
Геометрическим инвариантом называется величина или структура, сохраняющаяся при всех автоморфизмах реляционной структуры.
\end{definition}

\begin{proposition}
Геометрия системы определяется классом инвариантов её реляционной структуры.
\end{proposition}

\chapter{Выведение времени}

\section{Время как порядок изменений}

Время возникает не как внешний параметр, а как способ упорядочения различимых состояний.

\begin{definition}[История системы]
Историей системы \(\SubParts\) называется последовательность состояний
\[
\mathcal H=(\State_0,\State_1,\State_2,\dots),
\]
где соседние состояния различимы:
\[
\State_k\neq\State_{k+1}.
\]
\end{definition}

\begin{definition}[Реляционная длительность]
Длительность между двумя соседними состояниями определяется как
\[
\Delta t_k:=T(\State_k,\State_{k+1}).
\]
\end{definition}

\section{Отсутствие времени в статической системе}

\begin{theorem}[Статическая система атемпоральна]
Если история системы имеет вид
\[
\mathcal H=(\State,\State,\State,\dots),
\]
то реляционное время в ней не возникает.
\end{theorem}

\begin{proof}
Для любых соседних членов истории
\[
T(\State,\State)=0.
\]
Следовательно, отсутствует различимое изменение. Поэтому реляционное время не возникает.
\end{proof}

\section{Направление времени}

\begin{definition}[Ориентация истории]
Ориентацией истории называется выбор порядка
\[
\State_0\to\State_1\to\State_2\to\cdots.
\]
\end{definition}

\begin{remark}
Направление времени в ТЕЦ не является исходной абсолютной стрелой. Оно должно быть связано с асимметрией последовательности состояний или с необратимостью преобразований отношений.
\end{remark}

\chapter{Наблюдатели и относительные центры}

\section{Наблюдатель как часть}

\begin{definition}[Внутренний наблюдатель]
Внутренним наблюдателем называется часть \(O\in\Parts\), относительно которой строится описание отношений:
\[
\Observer_O:=\{\Rel(O,A_j)\mid A_j\in\Parts,\ A_j\neq O\}.
\]
\end{definition}

\begin{theorem}[Каждая часть может быть центром описания]
Для любой части \(A_i\in\Parts\) можно построить описание системы относительно неё.
\end{theorem}

\begin{proof}
Для произвольной части \(A_i\) определим
\[
\mathcal R_{A_i}:=\{\Rel(A_i,A_j)\mid A_j\in\Parts,\ A_j\neq A_i\}.
\]
Эта совокупность отношений задаёт описание системы относительно \(A_i\). Так как \(A_i\) выбрана произвольно, такое описание возможно для любой части.
\end{proof}

\section{Относительность описания}

Пусть \(A_i\) и \(A_k\) --- две разные части. Тогда возможны два описания:
\[
\mathcal R_{A_i}=\{\Rel(A_i,A_j)\}_{j\neq i},
\]
\[
\mathcal R_{A_k}=\{\Rel(A_k,A_j)\}_{j\neq k}.
\]

Эти описания не обязаны совпадать, но должны быть согласованы как разные представления одной и той же реляционной структуры.

\begin{definition}[Согласованность наблюдателей]
Два описания \(\mathcal R_{A_i}\) и \(\mathcal R_{A_k}\) называются согласованными, если существует преобразование \(\Phi_{ik}\), переводящее одно описание в другое без изменения инвариантов структуры:
\[
\Inv(\mathcal R_{A_i})=\Inv(\Phi_{ik}(\mathcal R_{A_i})).
\]
\end{definition}

\chapter{Реляционная механика}

\section{Движение}

\begin{definition}[Реляционное движение]
Движением части \(A_i\) относительно части \(A_j\) называется изменение отношения между ними:
\[
\Motion(A_i,A_j):=\Delta\Rel(A_i,A_j).
\]
\end{definition}

В этой формулировке движение не является перемещением в заранее данном пространстве. Движение есть изменение отношения.

\section{Скорость}

Если определены реляционное расстояние и реляционное время, можно ввести скорость:
\[
v(A_i,A_j):=\frac{\Delta d(A_i,A_j)}{\Delta t}.
\]

Подставляя определения, получаем
\[
v(A_i,A_j):=
\frac{\Delta F(\Rel(A_i,A_j))}
{\Delta T(\State_1,\State_2)}.
\]

\section{Ускорение}

Если скорость изменяется между состояниями, можно определить реляционное ускорение:
\[
a(A_i,A_j):=\frac{\Delta v(A_i,A_j)}{\Delta t}.
\]

\section{Инерция}

\begin{definition}[Реляционная инерция]
Инерцией системы называется устойчивость её отношений относительно изменений состояний.
\end{definition}

Формально можно ввести функционал устойчивости
\[
I(A_i):=\operatorname{Stab}\bigl(\{\Rel(A_i,A_j)\}_{j\neq i}\bigr),
\]
где \(\operatorname{Stab}\) измеряет сопротивление структуры отношений изменению.

\section{Масса и энергия как производные величины}

В ТЕЦ масса и энергия не вводятся исходно. Они должны быть выражены через устойчивость и изменение отношений.

Возможная схема:
\[
m(A_i):=M\bigl(\operatorname{Stab}(P(A_i))\bigr),
\]
где
\[
P(A_i)=\{\Rel(A_i,A_j)\mid j\neq i\}.
\]

Энергия может быть связана с мерой изменения состояния:
\[
E(\SubParts):=\mathcal E(\Delta\State),
\]
или более явно:
\[
E(\SubParts):=\mathcal E\bigl(\State_1(\SubParts),\State_2(\SubParts)\bigr).
\]

\begin{remark}
Эти формулы являются не завершёнными физическими законами, а направлением формализации. Их задача --- показать, что физические величины должны появляться как функции реляционной структуры.
\end{remark}

\chapter{Физическая интерпретация}

\section{Сопоставление с классической схемой}

Классическая схема описания мира:
\[
\text{пространство-время} \rightarrow \text{материя} \rightarrow \text{движение}.
\]

Схема ТЕЦ:
\[
\Whole \rightarrow \Parts \rightarrow \Rel \rightarrow \State \rightarrow \Space,\Time \rightarrow \text{физические величины}.
\]

Таким образом, пространство и время не являются сценой, на которой существуют отношения. Напротив, они являются устойчивыми формами организации отношений.

\section{Возможная связь с физикой}

ТЕЦ может быть сопоставлена с рядом идей современной физики:

\begin{enumerate}
    \item реляционными подходами к пространству и времени;
    \item идеей фоновой независимости;
    \item представлением о наблюдателе как внутренней части системы;
    \item геометрией как инвариантной структурой;
    \item временем как параметром изменения, а не внешней сущностью.
\end{enumerate}

Однако ТЕЦ не должна начинаться с готовых физических уравнений. Сначала требуется завершить внутреннюю математику теории.

\section{Критерии физической применимости}

Теория может получить физический смысл только при выполнении следующих условий:

\begin{enumerate}
    \item из отношений можно построить метрику или её обобщение;
    \item из изменения состояний можно получить устойчивую временную параметризацию;
    \item возможно определить наблюдателей и преобразования между их описаниями;
    \item известные физические величины выражаются как функции реляционных структур;
    \item теория даёт хотя бы одно проверяемое отличие от существующих моделей.
\end{enumerate}

\chapter{Открытые задачи}

\section{Задача о множестве значений отношений}

Необходимо определить природу множества \(\Values\). Возможны варианты:
\[
\Values=\{0,1\},
\]
\[
\Values\subseteq\mathbb R,
\]
\[
\Values \text{ является алгебраической структурой},
\]
\[
\Values \text{ является категорной или топологической структурой}.
\]

\section{Задача о метрике}

Нужно установить условия, при которых отображение
\[
F:\Values\to\DistVals
\]
порождает метрику.

\section{Задача о времени}

Нужно уточнить свойства функции
\[
T:\States\times\States\to\mathbb R_{\ge0}.
\]

В частности, важны вопросы:
\begin{enumerate}
    \item должна ли функция \(T\) быть симметричной;
    \item допускается ли направленное время;
    \item можно ли вывести необратимость;
    \item как связать \(T\) с физическим временем.
\end{enumerate}

\section{Задача о физических величинах}

Требуется строго определить:
\[
m,\quad E,\quad p,\quad F,\quad L,\quad H
\]
как производные от \(\Rel\), \(\State\), \(T\) и \(F\).

\chapter{Заключение}

Теория Единого Целого предлагает аксиоматический порядок, в котором пространство и время не являются фундаментальными сущностями. Исходными являются Единое Целое, части, отношения и состояния.

Главные выводы первой версии теории:

\begin{enumerate}
    \item существует только Единое Целое;
    \item всякая часть принадлежит Единому Целому и не равна ему;
    \item внешнего пространства и внешнего времени не существует;
    \item свойства частей возникают только через отношения;
    \item пространство возникает как структура отношений;
    \item время возникает как мера изменения отношений;
    \item движение есть изменение отношения;
    \item каждая часть может быть выбрана как относительный центр описания;
    \item физические величины должны быть выведены из реляционных структур.
\end{enumerate}

В краткой форме принцип ТЕЦ записывается так:
\[
\boxed{\Rel \prec \Space,\Time,\Physics.}
\]

Или словами:
\[
\boxed{\text{отношения первичны, пространство и время производны.}}
\]

\backmatter

\chapter*{План дальнейшей разработки}
\addcontentsline{toc}{chapter}{План дальнейшей разработки}

Дальнейшая монография может быть развита в следующих направлениях:

\begin{enumerate}[label=\textbf{Часть \Roman*.}, leftmargin=2cm]
    \item \textbf{Строгая теория отношений.} Классификация возможных \(\Values\), направленные и ненаправленные отношения, композиции отношений.
    \item \textbf{Реляционная топология.} Понятия близости, связности, окрестности и границы без предварительного пространства.
    \item \textbf{Реляционная геометрия.} Условия появления метрики, кривизны и размерности.
    \item \textbf{Реляционное время.} Истории состояний, необратимость, причинность и длительность.
    \item \textbf{Реляционная механика.} Движение, скорость, ускорение, инерция, масса и энергия.
    \item \textbf{Физическая проверка.} Сопоставление с классической механикой, общей теорией относительности и квантовыми реляционными подходами.
\end{enumerate}

\end{document}
